% Sección del artículo — dueño: P1.
\section{Conclusiones}
Presentamos NeuroFly-Bee, un controlador bioinspirado que separa el control local (mosca,
CNN+GRU) del retorno visual (abeja, red siamesa) y lo protege con un supervisor de seguridad.
En un simulador reproducible con seis condiciones, las líneas base clásicas completan como
máximo el \nfbMejorClasicoSinGPS{} de las misiones sin GPS y el \nfbMejorClasicoMoviles{} con
obstáculos móviles, y el análisis de errores atribuye esos fallos a la deriva y a planificar
sobre un mapa obsoleto. \pendiente{hallazgo principal del híbrido con los números de v19}

La validación del sensor muestra que la silueta 2.5D conserva las propiedades de los
panoramas reales que necesita el retorno visual, y la ablación en prueba mide cuánto aporta
comparar con la vista de casa: \nfbAblAbejaCompleta{} con ella frente a
\nfbAblAbejaVistaUnica{} sin ella.

Como trabajo futuro: (i) renderizar panoramas en 3D con dinámica de actitud, por ejemplo con
PyBullet~\cite{panerati2021}, aleatorizando texturas e iluminación~\cite{tobin2017};
(ii) evaluar la abeja sin reentrenar sobre las imágenes reales de bosque de Firlefyn
et~al.~\cite{firlefyn2024}, publicadas con licencia CC0, y reentrenar las arquitecturas de
Bee-Nav~\cite{ou2026} sobre nuestros datos para compararlas en igualdad de condiciones;
(iii) calibrar el viento con datos meteorológicos de cuencas peruanas y con el modelo de
turbulencia de Dryden; (iv) cuantizar las redes para un microcontrolador; y (v) validar el
sistema en un dron físico.
